\chapter{Revision history}
\label{app:revisions}

\section{History of versions}
\label{sec:version-history}

This is a theory under construction, and its structure and content have been remade several
times. The principal stages are recorded here. The names given to the versions are matters of
convenience.

The version number corresponds to the stage: stage $n$ is \texttt{v0.$n$.0}. A change of
structure, the addition or withdrawal of a proposition, or new empirical work raises the stage;
corrections of slips and adjustments of format are handled in the third digit (\texttt{v0.8.1},
for instance).

\begin{longtable}{@{}ll>{\raggedright\arraybackslash}p{0.18\textwidth}>{\raggedright\arraybackslash}p{0.48\textwidth}@{}}
\caption{History of versions: the version number, change of structure, and principal content of each stage.}
\label{tab:version-stages}\\
\toprule
Stage & Version & Structure & Principal content \\
\midrule
\endfirsthead
\toprule
Stage & Version & Structure & Principal content \\
\midrule
\endhead
\midrule
\multicolumn{4}{r@{}}{\footnotesize (continued on the next page)}\\
\endfoot
\bottomrule
\endlastfoot
Stage 1 & \texttt{v0.1.0} & chronological & the definition of $\Phi$, 27 types, the four-way classification of method \\
Stage 2 & \texttt{v0.2.0} & reordered logically & five parts: theory, catalogue, method, empirics, conclusion. Proofs and numerical examples added \\
Stage 3 & \texttt{v0.3.0} & change of parts & deductions under constraints moved from ``applications'' to ``the space of maps and constraints'' \\
Stage 4 & \texttt{v0.4.0} & empirics added & $\CCC$ and $g^\star$ measured from the Financial Statements Statistics \\
Stage 5 & \texttt{v0.5.0} & weaknesses classified & independence, exogeneity, and level versus change organized as three types \\
Stage 6 & \texttt{v0.6.0} & $\tau$ introduced & differences in level deduced from $\CCC=\sum\tau_i+\tau_{\mathrm{inv}}$ \\
Stage 7 & \texttt{v0.7.0} & existing theory integrated & existing theory matched to each term of $\kappa_i$; Appendix~\ref{app:priorwork} made three-layered \\
Stage 8 & \texttt{v0.8.0} & definitions filled in & the split of costs, $E$, $\kappa_K$, and the introduction of $d$ and $I$ into $g^\star$ \\
Stage 9 & \texttt{v0.9.0} & applied chapter added & deductions for replacing the supplier of the information structure (Chapter~\ref{ch:protocol}); the catalogue corrected to 28 entries \\
Stage 10 & \texttt{v0.10.0} & propositions filled in & simultaneity, identification of the beneficiary, and growth through prepayment added to Chapter~\ref{ch:protocol}; indivisibility imposed in Definition~\ref{def:protocol} \\
Stage 11 & \texttt{v0.11.0} & method filled in & revision requests recorded in the empirical chapters reflected into Chapter~\ref{ch:method}; failure of identification and the choice of variables added as sections \\
Stage 12 & \texttt{v0.12.0} & empirical design placed first & the five stages connecting theory to observation placed as Chapter~\ref{ch:design} at the head of Part~\ref{part:method}; the distinction of comparison from testing stated as a discipline \\
Stage 13 & \texttt{v0.13.0} & operationalization filled in & availability written as an equation and rules for assignment to types laid down; open problems reorganized into 21 items \\
Stage 14 & \texttt{v0.14.0} & measurement corrected & the effect of the scope of inventories on $\CCC$ and on the elasticity measured for ten industries (Section~\ref{sec:inventory-scope}) \\
Stage 15 & \texttt{v0.15.0} & further registrations & the correlation of size with turnover read as a proxy for B1, and implications Y and Z registered (Section~\ref{sec:b1-proxy}) \\
Stage 16 & \texttt{v0.16.0} & a route to testing Part~II & implications $\Phi$1--$\Phi$4 registered for the procedural records (Section~\ref{sec:protocol-predictions}) \\
Stage 17 & \texttt{v0.17.0} & implication Y tested & a national sample retrieved and tested on the part excluding Tokyo; partial support (Section~\ref{sec:predY-result}) \\
Stage 18 & \texttt{v0.18.0} & descriptions removed, format & a description resting on a single observation and an $n=2$ table removed; theorem-type body text unified in roman \\
Stage 19 & \texttt{v0.19.0} & means of reproduction included & data and processing placed in the repository, and the description in Appendix~\ref{app:data} brought into line with the fact \\
Stage 20 & \texttt{v0.20.0} & open problems reorganized & completed items dropped and errors of classification corrected, leaving 19 items \\
Stage 21 & \texttt{v0.21.0} & measurement extended, handling stated & the scope of inventories extended to 62 industries, weakening the verdict of stage 14 (Section~\ref{sec:inventory-scope}); the handling of per-firm data stated in Appendix~\ref{app:data} \\
Stage 22 & \texttt{v0.22.0} & grounds for selection, consistency of references & the reason for choosing families 2 and 5 derived from the three-way decomposition and the degrees of freedom (Section~\ref{ch:next}); a contradiction in the description of $\phi_{\rm cog}$ resolved (Section~\ref{sec:phicog-measurable}) \\
Stage 23 & \texttt{v0.23.0} & implication Z tested & 907 refund documents read, strength measured, and implication Z rejected (Section~\ref{sec:predZ-result}); registrations reached 18 and open problems 18 \\
Stage 24 & \texttt{v0.24.0} & implications $\Phi$1--$\Phi$4 tested & the deductions of Part~\ref{part:catalog} tested for the first time using the records of a public distributed ledger (Section~\ref{sec:protocol-results}); only $\Phi$2 supported; registrations reached 22 \\
Stage 25 & \texttt{v0.25.0} & quantities measured incidentally & $\lambda$ and the duration of procedures measured, giving an upper bound in Proposition~\ref{prop:unmanned-period} (Section~\ref{sec:protocol-byproducts}); implication D found identical to $\Phi$1 and its verdict changed to rejected \\
Stage 26 & \texttt{v0.26.0} & organized as a book & 22 chapters merged into 20 and the symbols of the implications cleared from the body; the wording of the registrations and their verdicts gathered into Appendix~\ref{app:prereg}, and the empirics of the procedure moved to the empirical part as Chapter~\ref{ch:protocol-emp}; verification obstacles presented as six classes from the outset \\
Stage 27 & \texttt{v0.27.2} & Part~I systematized & the chain of derivations starting from $\Phi$ carried through. Chapters 2 and 4 merged, reducing seven chapters to six; $\Delta$ and $\Pi$ defined from the primitives and $X=\Phi(\delta)$ adopted. The degrees of freedom of $\Phi$ organized into three groups and matched to the three ways of operating on risk \\
Stage 28 & \texttt{v0.28.0} & deductions gathered under the fundamental inequality & the cash constraint arranged as $\CCC\cdot r\leq E-A$ (Proposition~\ref{prop:fundamental}), from which the bound on scale, the time of depletion, the peak, the comparative statics and the switching of the binding constraint are derived. Antisymmetry and additivity of $\kappa$ added, and what the aggregate $\CCC$ measures identified. The funding constraint and the deadline of Chapter~\ref{ch:solo}, previously introduced separately, recast as special cases of the fundamental inequality \\
\end{longtable}


\section{Claims withdrawn}

Claims have been withdrawn several times in the course of this work. What was withdrawn and why is
recorded here. Since the body retains only the corrected description, which claims were once
thought to hold appears in this appendix alone.

So that a reader partway through is not misled, the body is kept consistent with the corrected
content. What follows is a record of how this work erred, and is not needed to understand the
body.

\section{Corrections deriving from measurement}

\begin{center}
\begin{tabularx}{\textwidth}{@{}XXX@{}}
\toprule
Original description & After correction & Occasion \\
\midrule
The rise in $\CCC$ lowered $g^\star$ by a third over a quarter-century
& $m$ rose 91\%, exceeding the 49\% deterioration of the denominator, and $m/\CCC$ rose from $38.5\%$ to $49.2\%$
& obtaining $m$ (Proposition~\ref{prop:gstar-flat}) \\
\addlinespace
The constraint on small firms is credit ($\kappa$)
& what actually binds is $m$; the $\CCC$ of small firms is shorter than that of large firms
& measurement by size stratum (Proposition~\ref{prop:small-margin}) \\
\addlinespace
$\phi_{\rm cog}$ has been measured nowhere in the world
& instances of measurement exist, and the method is established
& checking prior work (Section~\ref{sec:phicog-measurable}) \\
\addlinespace
The low $m$ of small firms is due to insufficient spreading of fixed costs
& officers' remuneration is the main cause; adding it back erases the ordering by size
& obtaining the breakdown of personnel costs (Proposition~\ref{prop:officer}) \\
\addlinespace
In R2, $\phi$ is demoted from the objective to a constraint
& withdrawn; the change is explicable by $H$ and $\beta$, and the demotion has no content
& checking prior work (Remark~\ref{rem:phi-demotion-retracted}) \\
\addlinespace
Time dependence of the capacity constraint should be taken into the theory
& withdrawn; it would rewrite the generating set of 27 types from an $n=1$ observation
& recognizing the asymmetry of the abstract and the concrete (Remark~\ref{rem:no-generalization}) \\
\addlinespace
The ill-posedness of implication B is a failure of the author's design
& it is the risk--incentive trade-off dispute, unresolved for a quarter-century
& checking prior work (Section~\ref{sec:predB-illposed}) \\
\bottomrule
\end{tabularx}
\captionof{table}{Corrections deriving from measurement: the original description, the corrected content, and the occasion.}
\label{tab:measurement-corrections}
\end{center}

The first of these is the weightiest. Having observed only the deterioration of $\CCC$, the work
came close to concluding that $g^\star$ had fallen; but Corollary~\ref{cor:gstar} is a ratio, and
no conclusion follows from a change in the denominator alone. The error would have been avoided by
obtaining $m$, and in fact the work had halted with the note that ``interpretation is reserved
because the change in $m$ has not been measured''. \textbf{Reserving judgement turned out to be
right.}

\section{Corrections deriving from reasoning}

\begin{center}
\begin{tabularx}{\textwidth}{@{}XX@{}}
\toprule
Original description & After correction \\
\midrule
Credit was ``being sold continuously'' during employment
& original attribution, as with works made for hire: it does not arise in the individual and then transfer \\
\addlinespace
The credit asset is greatest immediately after leaving employment
& what is greatest is not credit but the resource for conversion (contacts, technical knowledge); credit remains zero \\
\addlinespace
$\lambda_{\rm novel}$ is the sum of its components
& it is not additive: with a lock file, $\lambda_{\rm dep}$ is conditional on $\lambda_{\rm sec}$ \\
\addlinespace
Family 5 has favourable data conditions
& at the aggregate layer it is weaker than family 2; the claim holds only on descending to the per-firm layer \\
\addlinespace
Difficulty of access to real data is the principal obstacle to verification
& the principal obstacle for family 5 is (ii) ill-posedness of the hypothesis; family 2 is (iv) access constraint, but in the peculiar form that the required granularity is never published \\
\bottomrule
\end{tabularx}
\captionof{table}{Corrections deriving from reasoning: the original description and the corrected content.}
\label{tab:reasoning-corrections}
\end{center}

\section{A record of pre-registration at work}

The following are places where, had the implication not been fixed in writing beforehand, the work
would likely have drifted into retrospective interpretation.

\begin{enumerate}[label=(\arabic*)]
\item \textbf{The disappearance of capped-scheme fees.}
That 99.7\% of the capped scheme has disappeared in family 5 could have been read as
circumstantial support for the measurability hypothesis. In fact the price level of the capped
scheme is simply an order of magnitude lower, and it has nothing to do with measurability.

\item \textbf{The contemporaneous negative correlation of $\CCC$ with sales.}
A negative correlation was obtained in 31 of 45 industries ($p=0.008$). This could have been read
as ``$\CCC$ works as a business-cycle indicator'', but it is not distinguished from the
denominator effect.

\item \textbf{The credit position of large firms.}
That $\mathrm{DSO}-\mathrm{DPO}$ is large for large firms could have been interpreted as ``large
firms support their counterparties''. In fact it is the result of implication~F being rejected,
sign and all.
\end{enumerate}

The retrospective narrativization described in Section~\ref{sec:narrative} of others' accounts
arises equally in the author's own interpretation.

\section{Making assumptions explicit}

Several propositions here were originally stated without making their independence assumptions
explicit. In the course of verification it emerged that quantities in fact determined jointly were
being treated as independent, and the assumptions were written into each proposition.

\begin{center}
\begin{tabularx}{\textwidth}{@{}lXl@{}}
\toprule
Proposition & Assumption added & Occasion \\
\midrule
\ref{prop:pooling} & $\rho$ does not depend on $n$ & examining implication J \\
\ref{prop:breakage} & $\dbar$ is given exogenously & testing implication K (tiered menus) \\
\ref{prop:capacity} & $m$ and $\E[\delta]$ are independent & ibid. \\
\ref{prop:fcf} & $m$ and $\CCC$ are independent & obtaining $m$ (Proposition~\ref{prop:gstar-flat}) \\
\bottomrule
\end{tabularx}
\captionof{table}{Independence assumptions added to each proposition, and the occasion for adding them.}
\label{tab:assumption-additions}
\end{center}

The endogeneity of $\mathcal{G}$ (Remark~\ref{rem:G-endogenous}) was added at the same time. None
of the proofs was changed: \textbf{only the range of application was narrowed}.

A passage using the undefined symbol $\mathcal{X}$ as the domain of the production function was
rewritten without naming a domain. Since this work does not treat the internal structure of $f$,
no space of inputs need be introduced. $\mathcal{D}$, $\mathcal{F}$, $H$ and $\phi_{\rm risk}$
were added to the notation list.

\section{Introducing the lag}

The $\tau_i$ of Section~\ref{sec:tau} did not exist in the original version. It was introduced in
the course of working on weakness (C) organized in
Section~\ref{sec:level-change-gap}, the failure to separate level from change.

The decision to introduce it was made after first confirming that it would be used in several
places: the explanation of differences in level (Corollary~\ref{cor:ccc-level}), the impossibility
of cross-sectional comparison (Corollary~\ref{cor:no-cross-comparison}), and the decomposition of
change (Corollary~\ref{cor:ccc-change}).

Distortion of $\CCC$ during growth and a reinterpretation of the $\beta$ estimate were originally
also expected as uses, but the rejection of implications O and P showed that these cannot be
detected in annual data (Section~\ref{sec:lag-test}). That use was removed from the grounds for
introduction.

At the same time, Example~\ref{ex:insolvency} was corrected where it used $\CCC$ and $W$ before
their definitions, and rewritten in terms of $\kappa$ alone.

\section{Re-evaluating the relation to existing theory}

In the draft of this work, some of the terms of $\kappa_i$ were expected to be territory specific
to it, without checking whether existing theory covered them. On checking the literature, theory
existed for every one.

\begin{center}
\begin{tabularx}{\textwidth}{@{}lXl@{}}
\toprule
Original expectation & Existing theory found & Field \\
\midrule
Treating prepayment as trade credit is unexplored & reverse trade credit \cite{mateut2014,daripa2011} & corporate finance \\
There is no theory of $\kappa$ for deferred labour & the bonding argument \cite{akerlof1986} & labour economics \\
Accumulation of a public record is specific to this work & labour-market signalling \cite{spence1973} & economics of information \\
$\phi_{\rm cog}$ in family 2-4 is unmeasured & empirical work on breakage of prepaid balances \cite{nber2026pnbl} & consumer behaviour \\
\bottomrule
\end{tabularx}
\captionof{table}{For each term of $\kappa_i$: the original expectation, the existing theory found, and the field.}
\label{tab:prior-theory-check}
\end{center}

This check moved where the contribution lies. What was originally sought was \emph{territory
without existing theory}; in fact the contribution is \textbf{introducing existing theory as the
explanation of each term and treating them as a sum}
(Remark~\ref{rem:disjoint-literatures}). This is why Appendix~\ref{app:priorwork} was recast in
three layers.

The same failure has recurred. When implication B was set out in
Section~\ref{sec:predB-illposed}, the existence of the risk--incentive trade-off dispute had not
been checked.

\section{Filling in definitions}

The following are not corrections of errors but additions where the definitions were thin.

\begin{center}
\begin{tabularx}{\textwidth}{@{}lXl@{}}
\toprule
Subject & What was added & Occasion \\
\midrule
The technology side & costs split into $C_{\mathrm{prod}}+C_{\mathrm{admin}}$ (Definition~\ref{def:cost}) & the source of $\phi_{\mathrm{prod}}$ was undefined \\
$\phi_{\mathrm{prod}}$ & defined by marginal cost and a counterfactual price (Definition~\ref{def:phi-prod}) & ibid. \\
Equity $E$ & defined as paid-in capital plus accumulated $\phi$ (equation~\eqref{eq:equity}) & it was implicit as a residual \\
$\kappa_K$ & defined as a residual claim (Definition~\ref{def:kappa-K}) & equity was being treated identically to debt \\
$g^\star$ & recast to include depreciation $d$ and investment $I$ (Corollary~\ref{cor:gstar}) & $I\approx d\,r$ was implicitly assumed \\
\bottomrule
\end{tabularx}
\captionof{table}{What was added where definitions were thin, and the occasion.}
\label{tab:definition-supplements}
\end{center}

The revision of $g^\star$ changes the measured values. In the stratum below \textyen 10 million,
against a value of 20.9\% from $m$ alone, the nine-year average including $I$ is 3.1\%. The
conclusion of Proposition~\ref{prop:small-margin} is maintained, but the burden of investment is
added to the cause.

As to the split of costs, ``$C$ does not depend on $\Phi$'' was originally proved as a
proposition. But that $\Phi$ does not appear once $C(\delta)$ has been defined is a consequence of
the definition, and the proof was circular. It was withdrawn, rewritten as the definition of the
split, and the eight routes by which $\Phi$ generates costs enumerated.

Defining $\kappa_K$ revealed that subjectivity divides into three stages
(Remark~\ref{rem:subjectivity}). The subjectivity of the valuation map $v$ is asymmetry of
information; the subjectivity of the value $V_j$ of the residual claim is disagreement of
probability measures. The former can be resolved by disclosure; the latter cannot.

\section{Examining reversibility, and a correction}

The framework was reviewed from the standpoint of irreversibility and checked against existing
theory. One correction and two additions resulted.

\textbf{Correction.} In Remark~\ref{rem:rho-endogenous}, the mechanism making $\rho$ endogenous was
described as ``objects with low correlation are exhausted''. According to \cite{wagner2010} the
correct mechanism is different: it is not that the objects of pooling run out, but that
\textbf{the act of pooling itself creates correlation among those who pool}.
If two agents each hold half of two risks, each becomes individually safer while the two become
perfectly correlated.

\textbf{Additions.} The switching cost of $\Phi$ (Remark~\ref{rem:switching-cost}) and the social
optimality of pooling (Remark~\ref{rem:pooling-social}) were recorded. The former corresponds to
the theory of irreversible investment, the latter to research on systemic risk.

The framework of classifying reversibility by the size of a cost is itself already held by the
theory of irreversible investment and is not a contribution of this work. The formulation that a
decision is irreversible not because it cannot physically be undone but because undoing it is
costly belongs to the line of \cite{dixit1994}.

\section{Examination from the worker's viewpoint}
\label{sec:worker-predictions}

Whether to treat workers as an independent subject was examined, and the conclusion was not to.

The organization initially attempted was ``a worker is a special case of the one-person business'':
a one-person business whose $\Phi$ is given and whose customers concentrate in a single firm. It
emerged, however, that the constraints differ in both $\Phi$ and $\phi$, and this organization was
withdrawn (Remark~\ref{rem:worker-constraints}). The Labour Standards Act constrains five of the
six degrees of freedom of Section~\ref{sec:dof}, and $\phi_{\mathrm{prod}}$ and
$\phi_{\mathrm{cog}}$ have no counterparts.

Four implications were registered in the course of this examination.

\begin{center}
\begin{tabularx}{\textwidth}{@{}lXl@{}}
\toprule
Implication & Content & Result \\
\midrule
Q & workers with fewer assets have shorter payment cycles & untested \\
R & workers hold a higher $M_{\rm floor}$ than firms & untested \\
S & the paucity of types of employment contract is due to statutory constraint & untested \\
T & the new freelance statute shifts the allocation of $\kappa$ & \textbf{untestable} \\
U & the operating margin of individual enterprises is close to the corporate value with officers' remuneration added back & \textbf{supported} \\
\bottomrule
\end{tabularx}
\captionof{table}{Implications registered in the examination from the worker's viewpoint, and the results.}
\label{tab:worker-predictions}
\end{center}

Implication T is untestable because $\kappa$ for sole proprietors is not published. In this
process an obstacle was found that does not fit the four classes of
Section~\ref{sec:obstacle-taxonomy}: the form in which \textbf{a survey is taken but not
tabulated} (Remark~\ref{rem:tabulation-choice}).

\section{Extending the taxonomy of obstacles}

The four classes of Section~\ref{sec:obstacle-taxonomy} were constructed from the initial
experience of investigation. Later investigation produced two obstacles that do not fit them.

\begin{center}
\begin{tabularx}{\textwidth}{@{}lXl@{}}
\toprule
Class & Occasion of discovery & Record \\
\midrule
(v) choice of tabulation & seeking $\kappa$ for sole proprietors: the questionnaire exists but the tables do not & Section~\ref{sec:not-tabulated} \\
(vi) non-operation of the adjustment & in testing implication X, the regressor does not move although data for 307 firms exist & Remark~\ref{rem:no-adjustment} \\
\bottomrule
\end{tabularx}
\captionof{table}{The obstacles that did not fit the four classes: class, occasion of discovery, and where recorded.}
\label{tab:obstacle-extension}
\end{center}

(vi) in particular carries large implications for the method of this work. Part~\ref{part:method}
attributed the failure to verify to four causes: measurement, theory, identification and access.
The case in which \textbf{reality does not behave as the theory posits} had not been anticipated.

At the stage of setting out implication B for the measurability hypothesis, it was written that
$b$ could be obtained only as a binary. Later investigation revealed that the realized fee rate is
published as a continuous quantity, making the test possible. The obstacle was not the granularity
of the data.

\section{Extending the taxonomy of biases}

The selection biases of Part~\ref{part:method} originally consisted only of those concerning the
selection of the sample. Two distortions that the selection of the sample cannot explain emerged
in the course of the investigation.

\begin{center}
\begin{tabularx}{\textwidth}{@{}lXl@{}}
\toprule
Class & Occasion of discovery & Record \\
\midrule
Mistaking the unit & rows of a listing are offices while the values are replicated at firm level & Remark~\ref{rem:jinzai-dedup} \\
Selection of variables & in family 5, the fields of greatest theoretical interest are those whose aggregates cannot yield the required quantity & Remark~\ref{rem:data-gap-relevance} \\
\bottomrule
\end{tabularx}
\captionof{table}{Two distortions the selection of the sample cannot explain, and the occasions of discovery.}
\label{tab:bias-extension}
\end{center}

Both were added as sections of Chapter~\ref{ch:method}
(Sections~\ref{sec:misidentification} and \ref{sec:variable-selection}), and the opening of that
chapter was changed from ``six kinds'' to ``eight kinds''.

The former concerns \textbf{the unit of observation rather than the selection of the sample}, the
latter \textbf{variables rather than samples}. Both lay outside the range Chapter~\ref{ch:method}
originally treated.

\textbf{The grounds for the former were replaced at stage 17.} They originally rested on the
behaviour that, in retrieval specifying a licence number directly, a number that failed to resolve
returned the immediately preceding result. This admits interpretation as a measure limiting
automated retrieval, and erecting a general obstacle from a single observation was a leap of the
same form as in Section~\ref{sec:cap-observation}. The grounds were changed to the fact that in
Section~\ref{sec:predY-result} the presence or absence of consolidation reversed the sign of the
correlation. The description of that behaviour and the accompanying $n=2$ table were removed.

At the same time, the point that a complete frame does not yield power if the population is small
was added to Section~\ref{sec:frame}. Section~\ref{sec:nintei-small} had referred to that section
with ``as stated'', but the statement referred to did not exist there.

\section{Reorganizing the list of open problems}

The list, extended from 12 items to 21 at stage 13, was reorganized into 19 at stage 20.

\begin{enumerate}[label=(\arabic*)]
\item \textbf{A completed item remained.} Geographical extension of the sample was completed when
the national sample was retrieved for the test of Section~\ref{sec:predY-result}. It nevertheless
remained first among the items that could be started.

\item \textbf{One task was split in two.} Testing implication Z and collecting the refund rates and
periods that are its regressor were listed as separate items. A test and the data missing for it
are two faces of the same task; they were merged.

\item \textbf{A classification was wrong.} The curvature of the age effect was placed under ``items
with no obstacle'', but \textbf{pursuing it requires a sample carrying firm age}. The Financial
Statements Statistics used here has three axes — industry, size and year — and does not include
the year of founding. It was moved to (iv) access constraint.
\end{enumerate}

Items with no obstacle thereby fell from six to three. \textbf{A list drifts from the actual state
not when items are added but when finished items are left undeleted.}

\section{The means of reproduction were not included}

Appendix~\ref{app:data} stated that the processing used in the calculations was ``attached''.
\textbf{Not one file was in the repository.} A claim of the body was not satisfied by what was
published.

Data and processing were included at stage 19, and three defects were fixed at the same time.

\begin{enumerate}[label=(\arabic*)]
\item \textbf{The \texttt{parse.py} that Appendix~\ref{app:data} named did not run.} Its inputs
were still absolute paths from an old environment, and running it failed with files not found. It
was rewritten to read \texttt{raw/hojin/}.

\item \textbf{Some intermediate data had no means of regeneration.} No script existed to build the
series for officers' remuneration and depreciation, and they could not be restored from a fresh
clone. \texttt{parse.py} was extended to all six series, and the regenerated values were confirmed
to agree throughout with the previous ones.

\item A collision of output paths was resolved: the results of the test of implication E were
being written under the same name as intermediate data.
\end{enumerate}

After inclusion, it was confirmed from a state with \texttt{derived/} deleted that
\texttt{parse.py} and all 14 analyses run.

\section{Removing a description resting on a single observation}

Chapter~\ref{part:obstacles} recorded the behaviour that, in retrieval specifying a licence number
directly, a number that failed to resolve returned the immediately preceding result, and
generalized this into the obstacle ``retrieval succeeds but the object is mistaken''.

\textbf{It admits interpretation as a measure limiting automated retrieval, and other explanations
were not excluded.} Erecting a general obstacle from a single observation is a leap of the same
form as the one Remark~\ref{rem:no-generalization} warns against. The description and the
accompanying $n=2$ table were removed at stage 18.

The section of Chapter~\ref{ch:method} itself was retained. Its grounds were changed to the fact
that in Section~\ref{sec:predY-result} the presence or absence of consolidation reversed the sign
of the correlation, and its subject changed from ``mistaking'' to \textbf{the mismatch between what
a row denotes and what its values denote}. That appears at a magnitude sufficient to reverse the
conclusion in a sample of 1,402 firm rows across four occupations (1,130 distinct firms), and is
not a single observation.

\section{Setting theorem-type body text in roman}

Propositions and corollaries were being set in the \texttt{plain} style, which italicizes the body
text. Japanese has no italic form, so it stayed upright, and \textbf{only the Latin letters and
digits became italic, mixed in among it}. In places even the figures in tables placed inside
propositions were italicized.

They were unified in the same roman as definitions, remarks and examples. The headings remain bold
and numbered. Emphasis by \verb|\emph| in the body that contained Latin letters and became italic
was likewise changed to bold.

The heading of proofs was also changed. The default in amsthm is the English \texttt{Proof}, which
does not suit a Japanese document; \verb|\proofname| was set accordingly.

Italics thereby remain only in the punctuation attached to the proof heading and in journal names
in the bibliography.

\section{Testing implication Y, and a defect in the registration}

At stage 17 a national sample was retrieved from the Jinzai Service General Site and implication Y
tested on the part excluding Tokyo (Section~\ref{sec:predY-result}). The verdict is partial
support.

\textbf{The registration was defective.} Table~\ref{tab:predYZ-registration} laid down that offices
with amount-displayed fees would be used for implication Y, but the reference values of
Table~\ref{tab:b1-proxy-order} against which the comparison was made had been computed on a sample
excluding them (Remark~\ref{rem:predY-registration-error}). \textbf{The composition of the
reference values had not been checked at the time of registration.} Both rules are reported, and
the matter handled by showing that the verdict does not depend on the rule.

Three points emerged during retrieval.

First, \textbf{rows of the search results are offices while the values are at firm level}
(Remark~\ref{rem:jinzai-dedup}). Taking correlations without consolidating gives, among
physicians, a single company 134 times the weight, and the sign of the correlation turns from
positive to negative. Consolidation was not preprocessing but a procedure determining the
conclusion.

Second, \textbf{the displayed turnover rate cannot be computed from the placements and separations
in the same row}. For physicians the two agree in only 115 of 479 rows. The composition matches
Section~\ref{sec:predX-result} and does not impede comparison, but it was recorded that size and
turnover are not quantities over the same period (Remark~\ref{rem:predY-limits}).

Third, \textbf{the initial tabulation dropped part of the sample}. Extraction was restricted to the
licence-number letters Yu and Fu, so 200 Mu (free placement office) and 7 Toku records were
missed. This was wrongly described as loss during copying, and up to 11.2\% was reported as lost.
\textbf{The number of data rows agrees with the search results; nothing was lost.}

The error changed the verdict. With the dropped sample, the ordering excluding Tokyo did not agree
with the reference values, and it had been recorded that the ordering is not preserved. After
correction the ordering agrees perfectly under the registered rule. The verdict was set to partial
support not because the ordering is not preserved but because \textbf{the agreement of the
ordering depends on the rule of tabulation, and that rule had not been aligned with the
composition of the reference values at the time of registration}.

The reference values of Table~\ref{tab:b1-proxy-order} agree, on the Tokyo portion of the corrected
sample under the same composition, to within 0.02 in correlation; nurses agree to three decimal
places.

\section{A route for testing the deductions of Part~II}

The deductions of Part~\ref{part:catalog} \textbf{had almost no pre-registered implications} up to
stage 15. Of the 16 implications then registered, only D and T clearly targeted that part; D was
supported as to direction only, and T is untestable.

At stage 16, implications $\Phi$1--$\Phi$4 were registered for the ``inside'' partition of
Chapter~\ref{ch:protocol} (Section~\ref{sec:protocol-predictions}). In that partition the record of
the procedure is itself the population, and halted procedures remain, so the survival conditioning
of Chapter~\ref{ch:method} does not arise. The reason for abandonment given in
Chapter~\ref{ch:platformfee} — ``no population frame'' — disappears here.

Three points were distinguished before registering (Table~\ref{tab:protocol-testability}).

\begin{enumerate}[label=(\arabic*)]
\item \textbf{Quantities the procedure enforces are not objects of verification.} The collateral
constraint of Proposition~\ref{prop:collateral} is such a quantity: observing that collateral is
posted is merely reading the rules. Only quantities appearing as the result of an operator's choice
are objects of verification.

\item \textbf{The population differs} (Remark~\ref{rem:protocol-population}). What can be verified
is the correctness of the deduction, not whether the deduction applies to real one-person
businesses.

\item \textbf{Quantities at the level of agents cannot be recovered}
(Remark~\ref{rem:additivity}). What can be measured is only the distribution at the level of
identifiers.
\end{enumerate}

For $\Phi$1 it was noted that, because the number of dependencies may correlate with usage, the
alternative is hard to reject (Remark~\ref{rem:dep-usage-confound}).

\section{Registering implications Y and Z}

The observation of Remark~\ref{rem:scale-turnover} — that the correlation of size with turnover
differs greatly by occupation — \textbf{remained recorded as a secondary observation with no use}
up to stage 14.

At stage 15 a route was set out for reading it as a proxy for the B1 side, the informativeness of
effort that Table~\ref{tab:bstar-requirements} records as lacking a proxy, and implications Y and Z
were registered (Section~\ref{sec:b1-proxy}).

Two points were fixed at registration.

First, \textbf{the test must not use the same sample}
(Remark~\ref{rem:no-same-sample-test}). The correlation was observed after the fact in the 307-firm
sample, and testing on the same sample would be an error of the same form as retrospective
narrativization. What is required is reproduction out of sample.

Second, \textbf{the mechanism is not identified}
(Remark~\ref{rem:b1-proxy-assumption}). Thinning of input and looser screening give the same sign,
so the variable can be used as a proxy but no claim about the mechanism can be made.

There is a literature in labour economics on the relation between size and the quality of
placements which this work has not surveyed; it was added to ``theory that should be connected but
is not'' in Appendix~\ref{app:priorwork}.

\section{Checking the scope of inventories}

Chapter~\ref{part:industry} computed $\mathrm{DIO}$ from ``finished goods or merchandise'' alone,
excluding work in process and raw materials. The limitation was recorded as ``the effect on
analyses of differences and correlations is small, but the absolute level cannot be interpreted'',
\textbf{with no grounds given}.

At stage 14 the difference was measured against ``inventories (end of period)'' in the Financial
Statements Statistics (Section~\ref{sec:inventory-scope}). The result supports that statement: the
level moves by more than 30\% in some industries, but the correlation of $\Delta\CCC$ is 0.976.

Incidentally, the elasticity $\beta$ rose from 0.900 to 0.960 and the coefficient of determination
improved from 0.092 to 0.213. \textbf{Working capital is more nearly proportional to sales when
work in process and raw materials are included.} The contribution of the denominator effect moves
smaller, and the conclusion of Section~\ref{sec:denominator-effect} is strengthened.

During estimation it emerged that, because the observations with $W>0$ differ by definition, the
sign comes out reversed unless the sample is aligned (Remark~\ref{rem:common-sample}). Without
alignment the values are 0.894 and 0.787, giving the opposite conclusion that the denominator
effect increases.

\section{Filling in operationalization and reorganizing the list of open problems}

Classifying each item against the five stages of Chapter~\ref{ch:design} revealed that
\textbf{items halted at stage 1 (operationalization) were not registered in the list of open
problems}. Although a defect at stage 1 requires no data retrieval and is filled by definition, the
list was skewed towards problems of measurement and access.

The following were done at stage 13.

\begin{enumerate}[label=(\arabic*)]
\item \textbf{Availability was written as an equation} (Definition~\ref{def:availability}).
Previously it was stated only as ``consuming the constraint'' and had no form as an inequality.
Whereas the capacity constraint bounds the total volume of delivery, availability bounds the length
of an interval during which no response can be made. Both have the dimension of time.

\item \textbf{The reason the degree of automation cannot be defined was identified}
(Remark~\ref{rem:automation-rate-illposed}). It had been recorded that ``$\delta$ has no unit'',
but $\delta$ carries monetary units through the valuation map. Correctly, what $v$ measures is
value, not effort; and measuring on the response side leaves no measure defined on
$\Omega\setminus\hat\Omega$. \textbf{The obstacle is the absence of a measure, not of a unit, and
is not resolved by refining the definition.}

\item \textbf{Rules for assignment to types were laid down}
(Section~\ref{sec:type-assignment}). Using the catalogue empirically requires a rule for assigning
an observed $\Phi$ to a type, and none existed. Since the conditions identifying families are not
exclusive, an order was fixed. The choice of order has no deductive grounds and is a convention, so
it must be fixed in advance (Remark~\ref{rem:assignment-order-dependence}).

\item \textbf{A limitations section was placed in Chapter~\ref{part:unmanned}}
(Section~\ref{sec:unmanned-limits}). Of the deductive chapters it alone had none, and its
limitations were written in Chapter~\ref{ch:platformfee}. The theoretical limitations were moved to
that chapter and the empirical ones left in Chapter~\ref{ch:platformfee}.

\item \textbf{Open problems were reorganized from 12 items to 21.} Nine items in the limitations
sections and records of abandonment across the chapters were not reflected in the list. Six of them
\textbf{fell under no obstacle at all and had simply not been started}, so a subsection was
separated out for ``items with no obstacle''. The table of priorities was likewise split into what
can be started and what will not advance.
\end{enumerate}

This changed the close of Chapter~\ref{ch:conclusion}. It had read ``no item remains that can be
advanced by retrieving public data alone''; six items have no obstacle.

\section{Placing the empirical design first}

Chapter~\ref{part:obstacles} recorded that the four-way classification should be placed ahead of
the design of an investigation, because Part~\ref{part:method} treated selection bias in samples
while having no procedure for judging ``is verification possible at all?''.

This was carried out at stage 12 as Chapter~\ref{ch:design}. The four-way classification itself is
not placed first, however. The classification was extended to six through the empirical work, and
placing it first would put findings not yet obtained ahead of their derivation. \textbf{What was
brought forward is only the stages to be checked; the classification of causes for failing to pass
them was left in Chapter~\ref{part:obstacles}} (Remark~\ref{rem:taxonomy-placement}).

The distinction between comparison and testing was also stated as a discipline
(Section~\ref{sec:comparison-vs-test}). This work contains several consistency checks not resting
on pre-registration, such as the family-3 comparison of Chapter~\ref{ch:solo}, but no passage
stated their standing; each chapter merely noted individually that ``this is a retrospective
comparison and does not establish causation''. The discipline is needed because a comparison
requires no declaration of a population and so can be carried out while skipping stages 2 and 3
(Remark~\ref{rem:comparison-skips-stages}).

\section{Implication Z tested}

Up to stage 22, implication Z was left untested as ``the sample is already retrieved; what is
missing is only the strength of the refund scheme''. It was tested at stage 23 and
\textbf{rejected}.

\subsection{Implications \texorpdfstring{$\Phi$}{Phi}1--\texorpdfstring{$\Phi$}{Phi}4: three different failures}

Registered at stage 16 and tested at stage 24. Of the four, only $\Phi$2 is supported, and the
forms of failure all differ.

\textbf{$\Phi$1: an error in the design of the estimate.} At registration it was thought that
keeping the identifier as the unit of observation and using standard errors robust to clustering by
implementing address would handle the concentration of replicates. This is wrong. Robust standard
errors correct for within-cluster correlation but not for bias in the point estimate when one
cluster accounts for 95\% of the exposure. Only on reporting an estimate excluding the dominant
cluster did it emerge that the sign of the coefficient reverses and becomes insignificant
(Remark~\ref{rem:clone-concentration}).

\textbf{$\Phi$3: an error in the choice of control.} ``The distribution of $\CCC$ is skewed
non-negative'' was designed as a contrast with the corporate sector, but the corporate sector's
$\CCC$ is also non-negative throughout at the level of industry aggregates. Negative $\CCC$ appears
in individual firms and cancels out in aggregation. The implication and its alternative say the
same thing, and observation cannot tell them apart. This is ill-posedness of the same form as
implication B in Section~\ref{sec:predB-illposed} (Remark~\ref{rem:phi3-illposed}).

\textbf{$\Phi$4: mistaking the implication.} Proposition~\ref{prop:beneficiary} states that a
$\Phi$ requiring identification of the beneficiary cannot be implemented. From this ``only 7-4
survives among family 7'' was derived, but the proposition speaks of implementability, not
frequency of appearance. The quantitative part — that family 7 is rare — was right; the part about
the surviving type was wrong (Remark~\ref{rem:implementable-not-observed}).

\textbf{What the three have in common is that each was noticeable at the time of registration.}
All are problems of design that could have been examined before measurement; the data were not
lacking.

\subsection{Dropping the symbol for the production function}
\label{sec:drop-f}

Up to stage 26 the technology side was called a production function $f$. But \textbf{no functional
form was ever given}, and the notation list recorded that ``the domain is not stated''. There were
five mentions in the body, none of which used $f$ itself; it appeared only through
$C_{\mathrm{prod}}$.

At stage 27, $f$ was dropped and the technology side unified into $C_{\mathrm{prod}}$ (the
correspondence from delivery to cost). The definition of production surplus
(Definition~\ref{def:phi-prod}) already referred to $c(\delta)=dC_{\mathrm{prod}}/d\delta$ and
holds without $f$.

\textbf{A symbol with a name and no substance had been left in place.}

\subsection{Defining the relation between classes of agents and individuals}

The primitives merely listed customers $C$, suppliers $S$ and so on as ``the breakdown of
agents'', without giving the relation between an individual customer and the $\kappa_C$ denoting
customers as a whole. $\kappa_C$ occurs frequently in the body, and in
Example~\ref{ex:triple-index} the writing wavered between classes and individuals.

It was added to the definitions that a class is a subset of $N$, that an individual agent is
written as $i\in C$ and as an element of $N$, and that a class as a subscript denotes the sum
$\kappa_C=\sum_{i\in C}\kappa_i$. The index $j$ is used for identifiers in
Chapter~\ref{ch:protocol} and is therefore not used for individual agents.

\subsection{Revising the verdict on implication D}

Implication D was registered at stage 13 as ``the fewer the dependencies of a configuration, the
smaller $\lambda_{\rm dep}$ and the longer the period of unattended operation''. Implication $\Phi$1,
registered at stage 16, reads ``the fewer the other procedures a procedure calls, the longer the
period before it is repaired'' — \textbf{the same relation}. The duplication went unnoticed at
registration.

Up to stage 24 the table of verdicts recorded D as ``partial support''. But no section of the body
supported that verdict, and Chapter~\ref{ch:platformfee} recorded that ``implication D has not been
tested''. \textbf{The table and the body contradicted each other.}

The test of $\Phi$1 at stage 24 is also a test of D. The verdict was changed to \textbf{rejected}.

Registering the same relation twice is itself a failure in operating the pre-registration. The
cause was keeping the register by chapter and never reconciling it as a whole.

How strength was to be measured had not been laid down at registration. Deciding after reading the
documents would be fitting after the fact, so it was fixed as
equation~\eqref{eq:henreikin-strength} before they were read.

\textbf{Extraction by regular expression was abandoned.} Checking 8 documents produced errors in 3,
and the errors were biased towards overstating strength. Not being random error, they do not vanish
on averaging. All 907 were read by a person, and the 132 carrying no text layer were put through
OCR. None was illegible.

The reading brought several facts about publication practice to light. Some documents paste the
rate table as an image alone, so the body is legible while the content is not. Some offices have a
refund-field link pointing to a fee schedule. Some display the scheme as ``yes'' while filing a
rate of 0\%. \textbf{Filing a document and having its content legible are different things.}

\section{Giving deductive grounds for choosing families 2 and 5}

Section~\ref{ch:next} said only that ``families 2 and 5 contain the theoretical issues most
clearly'' and \textbf{did not show why those two}. What could be derived from the framework was
left underived.

It was rewritten in four steps. Since $\phi_{\mathrm{prod}}$ in equation~\eqref{eq:surplus} does
not depend on $\Phi$, the shape of $\Phi$ bears only on the remaining two terms.
$\phi_{\mathrm{cog}}$ requires degrees of freedom (5) and (6) and $\phi_{\mathrm{barg}}$ requires
(4), and the order of family assignment maps these onto families 2 and 5 respectively. The two
families sit on the diagonal of Figure~\ref{fig:quadrant}, spanned by the two dominant axes.

The reason for \emph{placing them side by side} was also set out first, as a difference of
institutional grounding. It does not anticipate the conclusion of Part~\ref{part:results}
(completeness and richness are independent) and is confined to the design-level reason that the two
need separating.

\section{Removing references to places that do not exist}

Part~\ref{part:results} contained ``at the end of Part~\ref{part:catalog} it was expected that
family 5 had the better conditions'', but \textbf{no such statement exists in
Part~\ref{part:catalog} or in Section~\ref{ch:next}}. A sentence from an earlier version had been
deleted, leaving only the sentence that answered it. The conclusion was kept and only the reference
to the premise dropped.

Also, the conditions for divergence were changed from a proposition to a remark at stage 2, but
\textbf{two places still referred to it as a proposition}. Both were corrected.

\section{Resolving a contradiction in the description of \texorpdfstring{$\phi_{\rm cog}$}{phi-cog}}

Immediately before Section~\ref{sec:phicog-measurable} it was stated that ``the only corresponding
measurement is the recovery figure for family 2'', while the section went on to say that ``prior
work contains direct instances of measurement''. \textbf{Of the same quantity it was written both
that no measurement exists and that measurements exist.}

Correctly, it is measurable and measurements exist; what this work obtained on its own is only an
upper bound. The family-2 recovery figure was folded in as one of the instances, and the passage
recast to list four.

\section{Extending the scope of inventories to 62 industries}

Stage 14 measured only ten industries and eleven size strata, and \textbf{recorded that ``the
elasticity improves on correcting the definition''}. Extending to 62 industries and all sizes
reverses the direction: the denominator effect rises slightly from 0.160 to 0.172.

The two are not error. Variation between size strata and variation between industries identify
different things. The estimate of that chapter is identified by variation between industries and
years, so it is the 62-industry figure that applies (Remark~\ref{rem:inventory-cut}).

\textbf{The range covered by the sample had been mistaken for the range reached by the conclusion.}
There was never any guarantee that a direction obtained on ten industries could be applied to a
45-industry estimate. All three verdicts remain ``supported'', but each is weaker than at stage 14.

\section{Stating the handling of per-firm data}

The Jinzai Service General Site publishes the filings of individual offices rather than aggregates.
These were included in the repository at stage 19, but \textbf{using them as input to analysis and
redistributing the collection are different things}.

The policy was written into Appendix~\ref{app:data} and the per-firm data removed from the
repository. What remains is only the list of licence and filing-receipt numbers, containing none of
the filed values. Re-retrieving with that as a seed reproduces the realized fee rates, turnover
rates and placement counts. What had already been published was removed from the history.

\section{Multiple definitions of a symbol}
\label{sec:rev-symbols}

At stage 28 the whole text was swept and six places where the same letter carried different
meanings were fixed.

\begin{center}
\begin{tabularx}{\textwidth}{@{}llX@{}}
\toprule
Letter & Meaning retained & Side changed \\
\midrule
$c$ & marginal cost $c(\delta)$ & flat fee $\to F$, living expenses $\to c_{\rm liv}$ \\
$m$ & gross or operating margin & number of members $\to n$ \\
$r$ & sales velocity & risk aversion $\to \gamma$ \\
$\tau$ & settlement lag $\tau_i$ & time of insolvency $\to \Tins$ \\
$A$ & fixed assets & the set of responses $\to \mathcal{A}$, agents $\to i,j$ \\
$a$ & response; age (APC) & allocation to commissioned work $\to \Ccl$ \\
\bottomrule
\end{tabularx}
\captionof{table}{Multiple definitions of symbols resolved at stage 28.}
\label{tab:symbol-collisions}
\end{center}

$c$ had four meanings. Only the cohort $c$ of age--period--cohort was retained, following the
notation of that literature, with a note to that effect in the body.

The three forms of degree of freedom (2) were also gathered into
equation~\eqref{eq:two-part}, $\pi=F+p\delta$: $p=0$ is flat, $F=0$ is usage-based, and both
non-zero is a two-part tariff. What had stood as three forms became one family.

\section{Correcting a count}

The catalogue of Part~\ref{part:catalog} was described consistently from stage 1 as ``27 types'',
but the counts by family are 4, 5, 3, 3, 5, 4, 4, totalling 28. Eleven places in the body were
corrected at stage 9. The composition of the families and the content of each type are unchanged.

\subsection{Calling the types a ``basis''}
\label{sec:rev-basis}

Up to stage 27 it was recorded of the catalogue of Chapter~\ref{ch:catalog} that ``the fact that a
$\Phi$ cannot be reduced to a single type means the types function as a basis''. But no addition
was defined on the space of $\Phi$, so the word had no meaning.

At stage 28 the additivity of $\kappa$ was laid down as
Proposition~\ref{prop:kappa-additive} and, at the same time, the non-uniqueness of the
decomposition shown (Remark~\ref{rem:decomposition-not-unique}). Without uniqueness the types are a
generating set, not a basis. The word in the body was changed.

\section{Considered but not adopted in the body}

In writing Chapter~\ref{ch:protocol} at stage 9, a structure presenting it as a novel contribution
was initially considered. On checking the literature it emerged that the restriction of funding by
a collateral constraint, the exclusion of agents holding no assets, and the divergence of
identifier from agent all have existing counterparts, so it was recast to claim no novelty and a
table of correspondences added to Appendix~\ref{app:priorwork}
(Remark~\ref{rem:collateral-priorwork}).

The chapter was retained because, in the position of showing how the theory of
Part~\ref{part:theory} functions under concrete conditions, an example varying the $\mathcal{G}$ of
Chapter~\ref{sec:info} was missing. Its role is the same as that of Chapter~\ref{ch:solo} and
Chapter~\ref{part:unmanned}.

\subsection{Describing an allocation that is not an interior solution as one}
\label{sec:rev-interior}

Up to stage 27 it was recorded of equation~\eqref{eq:leakage} in Chapter~\ref{ch:solo} that ``an
interior solution exists''. The equation is linear in the allocation $\Ccl$ and, under
$0<\theta<\gamma$, strictly decreasing, so taken alone it is maximized at the corner $\Ccl=0$. The
error arose from putting the cash that commissioned work brings in neither into the objective nor
into the constraints.

At stage 28 the cash the fundamental inequality demands was made explicit as a constraint, and the
optimal allocation re-derived as the edge of the constraint,
$\Ccl^\ast=c_{\rm liv}/w$ (Proposition~\ref{prop:allocation}). The direction of the conclusion —
that the criterion for selecting engagements moves from the fee to $\theta$ — is unchanged.

\section{Changes of structure}

Apart from corrections of content, the structure was changed twice.

\begin{enumerate}[label=(\arabic*)]
\item From the chronological order of the examination to the logical order of theory, catalogue,
method, empirics and conclusion.
\item The deductions under constraints (Chapters~\ref{ch:solo} and \ref{part:unmanned}) moved from
a separate part on ``applications'' into Part~\ref{part:catalog}. Applications are justified only
after a theory has passed through empirical work, and the empirical work here is not complete.
\end{enumerate}

The second change was accompanied by fixing 64 inconsistent references (parts confused with
chapters) and two logical inversions (a later chapter referred to in the past tense).
